% =====================================================================
%  ESEMPIO di documento diviso in capitoli: mostra tutto quello che il
%  tipo [documento] sa fare. Il contenuto è di prova.
%
%  Struttura della cartella:
%     esempio-documento.tex      <- questo file (il "principale")
%     capitoli/introduzione.tex  <- richiamato con % !TEX root = ../esempio-documento.tex
% (la riga qui sopra dice a VS Code quale file compilare quando salvate questo)
% Capitolo richiamato da esempio-documento.tex con % !TEX root = ../esempio-documento.tex
% (la riga qui sopra dice a VS Code quale file compilare quando salvate questo)
% Capitolo richiamato da esempio-documento.tex con % !TEX root = ../esempio-documento.tex
% (la riga qui sopra dice a VS Code quale file compilare quando salvate questo)
% Capitolo richiamato da esempio-documento.tex con \input{capitoli/introduzione}
% (niente \documentclass e niente \begin{document}: solo il testo)

\section{Introduzione}

\subsection{Scopo del documento}
Questo documento raccoglie le regole che il gruppo segue per lavorare:
come si scrivono i documenti, come si usa il repository e come si
verifica il lavoro degli altri.

\subsection{Glossario}
I termini che potrebbero non essere chiari sono segnati con una
\glo{G} a pedice, per esempio \glo{repository} o \glo{pull request},
e sono spiegati nel Glossario.

\subsection{Riferimenti}
\begin{itemize}
  \item Organizzazione GitHub del gruppo:
        \url{https://github.com/Git-Push-Team}
  \item Sito con i documenti del gruppo:
        \url{https://git-push-team.github.io/docs/}
\end{itemize}

% (niente \documentclass e niente \begin{document}: solo il testo)

\section{Introduzione}

\subsection{Scopo del documento}
Questo documento raccoglie le regole che il gruppo segue per lavorare:
come si scrivono i documenti, come si usa il repository e come si
verifica il lavoro degli altri.

\subsection{Glossario}
I termini che potrebbero non essere chiari sono segnati con una
\glo{G} a pedice, per esempio \glo{repository} o \glo{pull request},
e sono spiegati nel Glossario.

\subsection{Riferimenti}
\begin{itemize}
  \item Organizzazione GitHub del gruppo:
        \url{https://github.com/Git-Push-Team}
  \item Sito con i documenti del gruppo:
        \url{https://git-push-team.github.io/docs/}
\end{itemize}

% (niente \documentclass e niente \begin{document}: solo il testo)

\section{Introduzione}

\subsection{Scopo del documento}
Questo documento raccoglie le regole che il gruppo segue per lavorare:
come si scrivono i documenti, come si usa il repository e come si
verifica il lavoro degli altri.

\subsection{Glossario}
I termini che potrebbero non essere chiari sono segnati con una
\glo{G} a pedice, per esempio \glo{repository} o \glo{pull request},
e sono spiegati nel Glossario.

\subsection{Riferimenti}
\begin{itemize}
  \item Organizzazione GitHub del gruppo:
        \url{https://github.com/Git-Push-Team}
  \item Sito con i documenti del gruppo:
        \url{https://git-push-team.github.io/docs/}
\end{itemize}

%     capitoli/processi.tex      <- richiamato con % !TEX root = ../esempio-documento.tex
% (la riga qui sopra dice a VS Code quale file compilare quando salvate questo)
% Capitolo richiamato da esempio-documento.tex con % !TEX root = ../esempio-documento.tex
% (la riga qui sopra dice a VS Code quale file compilare quando salvate questo)
% Capitolo richiamato da esempio-documento.tex con % !TEX root = ../esempio-documento.tex
% (la riga qui sopra dice a VS Code quale file compilare quando salvate questo)
% Capitolo richiamato da esempio-documento.tex con \input{capitoli/processi}

\section{Processi di supporto}

\subsection{Documentazione}
Ogni documento ha un redattore e un verificatore diversi. Il ciclo di
vita di un documento è:
\begin{enumerate}
  \item stesura nel proprio branch;
  \item apertura della pull request;
  \item verifica sul PDF dell'anteprima automatica;
  \item approvazione e unione in \glo{main}.
\end{enumerate}

\begin{nota}[Da ricordare]
La versione del documento è quella della prima riga del registro delle
modifiche: non va scritta da nessun'altra parte.
\end{nota}

\subsection{Versioni}
La Tabella~\ref{tab:versioni} riassume quando cambia ogni numero della
versione.

\begin{table}[H]
  \centering
  \caption{Regole per il numero di versione\label{tab:versioni}}
  \begin{tabular}{@{}ll@{}}
    \toprule
    \intestazione{Quando} & \intestazione{Nuova versione} \\
    \midrule
    Prima stesura                & 0.1.0 \\
    Ogni modifica prima della verifica & 0.2.0, 0.3.0, \ldots \\
    Verifica e approvazione      & 1.0.0 \\
    Modifica dopo l'approvazione & 1.1.0, 1.2.0, \ldots \\
    \bottomrule
  \end{tabular}
\end{table}

\subsection{Architettura di riferimento}
Le immagini stanno nella cartella \emph{allegati} e si inseriscono con
\emph{figura}, come in Figura~\ref{fig:architettura}.

\figura[0.9]{esempi/2026-10-15/architettura-aura.png}{Un'immagine di
esempio presa dalla cartella allegati\label{fig:architettura}}

\section{Verifica}
\subsection{Controlli sul documento}
Il Verificatore legge il PDF per intero e controlla ortografia, date,
nomi e riferimenti, poi approva la pull request.


\section{Processi di supporto}

\subsection{Documentazione}
Ogni documento ha un redattore e un verificatore diversi. Il ciclo di
vita di un documento è:
\begin{enumerate}
  \item stesura nel proprio branch;
  \item apertura della pull request;
  \item verifica sul PDF dell'anteprima automatica;
  \item approvazione e unione in \glo{main}.
\end{enumerate}

\begin{nota}[Da ricordare]
La versione del documento è quella della prima riga del registro delle
modifiche: non va scritta da nessun'altra parte.
\end{nota}

\subsection{Versioni}
La Tabella~\ref{tab:versioni} riassume quando cambia ogni numero della
versione.

\begin{table}[H]
  \centering
  \caption{Regole per il numero di versione\label{tab:versioni}}
  \begin{tabular}{@{}ll@{}}
    \toprule
    \intestazione{Quando} & \intestazione{Nuova versione} \\
    \midrule
    Prima stesura                & 0.1.0 \\
    Ogni modifica prima della verifica & 0.2.0, 0.3.0, \ldots \\
    Verifica e approvazione      & 1.0.0 \\
    Modifica dopo l'approvazione & 1.1.0, 1.2.0, \ldots \\
    \bottomrule
  \end{tabular}
\end{table}

\subsection{Architettura di riferimento}
Le immagini stanno nella cartella \emph{allegati} e si inseriscono con
\emph{figura}, come in Figura~\ref{fig:architettura}.

\figura[0.9]{esempi/2026-10-15/architettura-aura.png}{Un'immagine di
esempio presa dalla cartella allegati\label{fig:architettura}}

\section{Verifica}
\subsection{Controlli sul documento}
Il Verificatore legge il PDF per intero e controlla ortografia, date,
nomi e riferimenti, poi approva la pull request.


\section{Processi di supporto}

\subsection{Documentazione}
Ogni documento ha un redattore e un verificatore diversi. Il ciclo di
vita di un documento è:
\begin{enumerate}
  \item stesura nel proprio branch;
  \item apertura della pull request;
  \item verifica sul PDF dell'anteprima automatica;
  \item approvazione e unione in \glo{main}.
\end{enumerate}

\begin{nota}[Da ricordare]
La versione del documento è quella della prima riga del registro delle
modifiche: non va scritta da nessun'altra parte.
\end{nota}

\subsection{Versioni}
La Tabella~\ref{tab:versioni} riassume quando cambia ogni numero della
versione.

\begin{table}[H]
  \centering
  \caption{Regole per il numero di versione\label{tab:versioni}}
  \begin{tabular}{@{}ll@{}}
    \toprule
    \intestazione{Quando} & \intestazione{Nuova versione} \\
    \midrule
    Prima stesura                & 0.1.0 \\
    Ogni modifica prima della verifica & 0.2.0, 0.3.0, \ldots \\
    Verifica e approvazione      & 1.0.0 \\
    Modifica dopo l'approvazione & 1.1.0, 1.2.0, \ldots \\
    \bottomrule
  \end{tabular}
\end{table}

\subsection{Architettura di riferimento}
Le immagini stanno nella cartella \emph{allegati} e si inseriscono con
\emph{figura}, come in Figura~\ref{fig:architettura}.

\figura[0.9]{esempi/2026-10-15/architettura-aura.png}{Un'immagine di
esempio presa dalla cartella allegati\label{fig:architettura}}

\section{Verifica}
\subsection{Controlli sul documento}
Il Verificatore legge il PDF per intero e controlla ortografia, date,
nomi e riferimenti, poi approva la pull request.

%  I file dei capitoli NON hanno \documentclass né \begin{document}:
%  contengono solo testo. Si compila sempre il file principale.
% =====================================================================
\documentclass[documento]{gpt}

\titolo{Norme di Progetto}
\uso{esterno}
\responsabile{spagnoli}

% Registro: la riga più recente in cima. In copertina compaiono da soli
% versione (1.1.0), data dell'ultima modifica, redattori e verificatori.
\modifica{1.1.0}{2026-11-20}{mazzurana}{cescobolla}{Aggiunta la sezione sulla verifica}
\modifica{1.0.0}{2026-11-10}{guglielmo}{corro}{Verifica e approvazione}
\modifica{0.2.0}{2026-11-05}{biasini}{--}{Stesura della sezione sui processi}
\modifica{0.1.0}{2026-10-28}{guglielmo}{--}{Prima stesura: struttura del documento}

\begin{document}

% !TEX root = ../esempio-documento.tex
% (la riga qui sopra dice a VS Code quale file compilare quando salvate questo)
% Capitolo richiamato da esempio-documento.tex con % !TEX root = ../esempio-documento.tex
% (la riga qui sopra dice a VS Code quale file compilare quando salvate questo)
% Capitolo richiamato da esempio-documento.tex con % !TEX root = ../esempio-documento.tex
% (la riga qui sopra dice a VS Code quale file compilare quando salvate questo)
% Capitolo richiamato da esempio-documento.tex con \input{capitoli/introduzione}
% (niente \documentclass e niente \begin{document}: solo il testo)

\section{Introduzione}

\subsection{Scopo del documento}
Questo documento raccoglie le regole che il gruppo segue per lavorare:
come si scrivono i documenti, come si usa il repository e come si
verifica il lavoro degli altri.

\subsection{Glossario}
I termini che potrebbero non essere chiari sono segnati con una
\glo{G} a pedice, per esempio \glo{repository} o \glo{pull request},
e sono spiegati nel Glossario.

\subsection{Riferimenti}
\begin{itemize}
  \item Organizzazione GitHub del gruppo:
        \url{https://github.com/Git-Push-Team}
  \item Sito con i documenti del gruppo:
        \url{https://git-push-team.github.io/docs/}
\end{itemize}

% (niente \documentclass e niente \begin{document}: solo il testo)

\section{Introduzione}

\subsection{Scopo del documento}
Questo documento raccoglie le regole che il gruppo segue per lavorare:
come si scrivono i documenti, come si usa il repository e come si
verifica il lavoro degli altri.

\subsection{Glossario}
I termini che potrebbero non essere chiari sono segnati con una
\glo{G} a pedice, per esempio \glo{repository} o \glo{pull request},
e sono spiegati nel Glossario.

\subsection{Riferimenti}
\begin{itemize}
  \item Organizzazione GitHub del gruppo:
        \url{https://github.com/Git-Push-Team}
  \item Sito con i documenti del gruppo:
        \url{https://git-push-team.github.io/docs/}
\end{itemize}

% (niente \documentclass e niente \begin{document}: solo il testo)

\section{Introduzione}

\subsection{Scopo del documento}
Questo documento raccoglie le regole che il gruppo segue per lavorare:
come si scrivono i documenti, come si usa il repository e come si
verifica il lavoro degli altri.

\subsection{Glossario}
I termini che potrebbero non essere chiari sono segnati con una
\glo{G} a pedice, per esempio \glo{repository} o \glo{pull request},
e sono spiegati nel Glossario.

\subsection{Riferimenti}
\begin{itemize}
  \item Organizzazione GitHub del gruppo:
        \url{https://github.com/Git-Push-Team}
  \item Sito con i documenti del gruppo:
        \url{https://git-push-team.github.io/docs/}
\end{itemize}

% !TEX root = ../esempio-documento.tex
% (la riga qui sopra dice a VS Code quale file compilare quando salvate questo)
% Capitolo richiamato da esempio-documento.tex con % !TEX root = ../esempio-documento.tex
% (la riga qui sopra dice a VS Code quale file compilare quando salvate questo)
% Capitolo richiamato da esempio-documento.tex con % !TEX root = ../esempio-documento.tex
% (la riga qui sopra dice a VS Code quale file compilare quando salvate questo)
% Capitolo richiamato da esempio-documento.tex con \input{capitoli/processi}

\section{Processi di supporto}

\subsection{Documentazione}
Ogni documento ha un redattore e un verificatore diversi. Il ciclo di
vita di un documento è:
\begin{enumerate}
  \item stesura nel proprio branch;
  \item apertura della pull request;
  \item verifica sul PDF dell'anteprima automatica;
  \item approvazione e unione in \glo{main}.
\end{enumerate}

\begin{nota}[Da ricordare]
La versione del documento è quella della prima riga del registro delle
modifiche: non va scritta da nessun'altra parte.
\end{nota}

\subsection{Versioni}
La Tabella~\ref{tab:versioni} riassume quando cambia ogni numero della
versione.

\begin{table}[H]
  \centering
  \caption{Regole per il numero di versione\label{tab:versioni}}
  \begin{tabular}{@{}ll@{}}
    \toprule
    \intestazione{Quando} & \intestazione{Nuova versione} \\
    \midrule
    Prima stesura                & 0.1.0 \\
    Ogni modifica prima della verifica & 0.2.0, 0.3.0, \ldots \\
    Verifica e approvazione      & 1.0.0 \\
    Modifica dopo l'approvazione & 1.1.0, 1.2.0, \ldots \\
    \bottomrule
  \end{tabular}
\end{table}

\subsection{Architettura di riferimento}
Le immagini stanno nella cartella \emph{allegati} e si inseriscono con
\emph{figura}, come in Figura~\ref{fig:architettura}.

\figura[0.9]{esempi/2026-10-15/architettura-aura.png}{Un'immagine di
esempio presa dalla cartella allegati\label{fig:architettura}}

\section{Verifica}
\subsection{Controlli sul documento}
Il Verificatore legge il PDF per intero e controlla ortografia, date,
nomi e riferimenti, poi approva la pull request.


\section{Processi di supporto}

\subsection{Documentazione}
Ogni documento ha un redattore e un verificatore diversi. Il ciclo di
vita di un documento è:
\begin{enumerate}
  \item stesura nel proprio branch;
  \item apertura della pull request;
  \item verifica sul PDF dell'anteprima automatica;
  \item approvazione e unione in \glo{main}.
\end{enumerate}

\begin{nota}[Da ricordare]
La versione del documento è quella della prima riga del registro delle
modifiche: non va scritta da nessun'altra parte.
\end{nota}

\subsection{Versioni}
La Tabella~\ref{tab:versioni} riassume quando cambia ogni numero della
versione.

\begin{table}[H]
  \centering
  \caption{Regole per il numero di versione\label{tab:versioni}}
  \begin{tabular}{@{}ll@{}}
    \toprule
    \intestazione{Quando} & \intestazione{Nuova versione} \\
    \midrule
    Prima stesura                & 0.1.0 \\
    Ogni modifica prima della verifica & 0.2.0, 0.3.0, \ldots \\
    Verifica e approvazione      & 1.0.0 \\
    Modifica dopo l'approvazione & 1.1.0, 1.2.0, \ldots \\
    \bottomrule
  \end{tabular}
\end{table}

\subsection{Architettura di riferimento}
Le immagini stanno nella cartella \emph{allegati} e si inseriscono con
\emph{figura}, come in Figura~\ref{fig:architettura}.

\figura[0.9]{esempi/2026-10-15/architettura-aura.png}{Un'immagine di
esempio presa dalla cartella allegati\label{fig:architettura}}

\section{Verifica}
\subsection{Controlli sul documento}
Il Verificatore legge il PDF per intero e controlla ortografia, date,
nomi e riferimenti, poi approva la pull request.


\section{Processi di supporto}

\subsection{Documentazione}
Ogni documento ha un redattore e un verificatore diversi. Il ciclo di
vita di un documento è:
\begin{enumerate}
  \item stesura nel proprio branch;
  \item apertura della pull request;
  \item verifica sul PDF dell'anteprima automatica;
  \item approvazione e unione in \glo{main}.
\end{enumerate}

\begin{nota}[Da ricordare]
La versione del documento è quella della prima riga del registro delle
modifiche: non va scritta da nessun'altra parte.
\end{nota}

\subsection{Versioni}
La Tabella~\ref{tab:versioni} riassume quando cambia ogni numero della
versione.

\begin{table}[H]
  \centering
  \caption{Regole per il numero di versione\label{tab:versioni}}
  \begin{tabular}{@{}ll@{}}
    \toprule
    \intestazione{Quando} & \intestazione{Nuova versione} \\
    \midrule
    Prima stesura                & 0.1.0 \\
    Ogni modifica prima della verifica & 0.2.0, 0.3.0, \ldots \\
    Verifica e approvazione      & 1.0.0 \\
    Modifica dopo l'approvazione & 1.1.0, 1.2.0, \ldots \\
    \bottomrule
  \end{tabular}
\end{table}

\subsection{Architettura di riferimento}
Le immagini stanno nella cartella \emph{allegati} e si inseriscono con
\emph{figura}, come in Figura~\ref{fig:architettura}.

\figura[0.9]{esempi/2026-10-15/architettura-aura.png}{Un'immagine di
esempio presa dalla cartella allegati\label{fig:architettura}}

\section{Verifica}
\subsection{Controlli sul documento}
Il Verificatore legge il PDF per intero e controlla ortografia, date,
nomi e riferimenti, poi approva la pull request.


\end{document}
